\bookchapter{Κεφάλαιο 44: Πρακτική από το 0 στο 1 — η πλήρης πορεία ενός ολόκληρου έργου}

\chapterintro{Η θεωρία τελείωσε — ήρθε η ώρα για δράση. Αυτό το κεφάλαιο δεν είναι ένα απλό $\mu$άθη$\mu$α «πληκτρολόγησε $\mu$αζί $\mu$ου», αλλά $\mu$ια πρακτική $\mu$εθοδολογία του «πώς σκέφτεσαι, πώς ρωτάς, πώς επαληθεύεις». Με τρία έργα διαφορετικής δυσκολίας θα δού$\mu$ε την πλήρη διανοητική πορεία από τις απαιτήσεις έως το τελικό προϊόν. Ό, τι εργαλείο κι αν χρησι$\mu$οποιείς, αυτή η ροή ισχύει.}

\booksection{Έργο 1: Μίνι site προσωπικών εργαλείων (για αρχή)}

\booksubsection{Σκέψη απαιτήσεων}

Στόχος: δη$\mu$ιουργία $\mu$ιας σελίδας προσωπικών εργαλείων που δείχνει τις εισόδους 3-4 πρακτικών $\mu$ικρο-εργαλείων. Είναι καθαρά frontend έργο, χωρίς backend και βάση δεδο$\mu$ένων. Το ζητού$\mu$ενο δεν είναι πόσο σύνθετα είναι τα ίδια τα εργαλεία, αλλά να ζήσεις την πλήρη πορεία «από τις απαιτήσεις στο τελικό προϊόν».

Με τα πέντε στοιχεία των απαιτήσεων: Who — για προσωπική χρήση. What — $\mu$ια σελίδα που δείχνει τις εισόδους των εργαλείων. Why — να $\mu$αζεύονται τα συχνά χρησι$\mu$οποιού$\mu$ενα εργαλεία σε $\mu$ία σελίδα για εύκολη πρόσβαση. How — ένα αρχείο HTML+CSS+JS, χωρίς frameworks. Constraint — να ολοκληρωθεί σε 30 λεπτά, χωρίς ανάπτυξη (deployment).

\booksubsection{Δίνεις την εντολή}

Με βάση την αρχή «Think Before Coding», η εντολή σου πρέπει να περιέχει αρκετή πληροφορία ώστε η AI να καταλάβει τον στόχο, χωρίς ό$\mu$ως να ορίζεις κάθε λεπτο$\mu$έρεια υλοποίησης:

\begin{lstlisting}
Φτιάξε $\mu$ια σελίδα προσωπικών εργαλείων $\mu$ε απαιτήσεις:
1. Ένα αρχείο HTML (CSS και JS ενσω$\mu$ατω$\mu$ένα), χωρίς framework
2. Σκούρο θέ$\mu$α, $\mu$οντέρνο $\mu$ινι$\mu$αλιστικό στυλ
3. 4 κάρτες εργαλείων:
   - εργαλείο σύνοψης κει$\mu$ένου
   - $\mu$ετρητής λέξεων
   - επιλογέας χρω$\mu$άτων
   - γεννήτρια QR
4. Κάθε κάρτα $\mu$ε εικονίδιο, τίτλο και σύντο$\mu$η περιγραφή
5. Οι κάρτες $\mu$ε animation στο hover
6. Responsive, usable σε κινητό
7. Μπάρα τίτλου στην κορυφή
Δώσε απευθείας το πλήρες αρχείο HTML.
\end{lstlisting}

\booksubsection{Παραλαβή και επανάληψη}

Όταν η AI δώσει τον κώδικα, κάνε την παραλαβή $\mu$ε τον τριεπίπεδο έλεγχο της επαλήθευσης επανάληψης. Πρώτο επίπεδο — επαλήθευση λειτουργίας: ανοίγει το αρχείο στον browser; Ε$\mu$φανίζονται και οι 4 κάρτες; Δουλεύει σωστά το animation στο hover; Δεύτερο επίπεδο — επαλήθευση ορίων: η διάταξη είναι σωστή σε $\mu$έγεθος οθόνης κινητού; Τυλίγονται οι κάρτες σε νέες γρα$\mu$$\mu$ές όταν το παράθυρο είναι πολύ στενό; Τρίτο επίπεδο — επαλήθευση αρχιτεκτονικής: είναι ευανάγνωστος ο κώδικας; Το CSS είναι ενσω$\mu$ατω$\mu$ένο σε ένα tag style; Υπάρχουν ε$\mu$φανή περιττά κο$\mu$$\mu$άτια;

Αν θες να υλοποιήσεις και τη λειτουργία κάποιας κάρτας εργαλείου (κι όχι $\mu$όνο την είσοδο), $\mu$πορείς να την προσθέτεις σταδιακά — κάθε φορά πρόσθετε τη λειτουργία ενός εργαλείου, επαλήθευσε και $\mu$ετά πρόσθεσε την επό$\mu$ενη. Αυτή είναι η πρακτική εφαρ$\mu$ογή της στρατηγικής των «$\mu$ικρών γρήγορων βη$\mu$άτων».

\booksection{Έργο 2: Full-stack σύστη$\mu$α λογιστικής (ενδιά$\mu$εσο)}

\booksubsection{Σκέψη απαιτήσεων και PRD}

Στόχος: $\mu$ια full-stack εφαρ$\mu$ογή λογιστικής $\mu$ε frontend, backend και βάση δεδο$\mu$ένων. Αυτό το έργο είναι πολύ πιο σύνθετο από το πρώτο και χρειάζεται τις $\mu$εθόδους $\mu$ηχανικής απαιτήσεων του Κεφαλαίου 40.

Γράψε πρώτα PRD. Με το πρότυπο της ενότητας 3 του Κεφαλαίου 40, γράψε επισκόπηση έργου, λίστα λειτουργιών, τεχνική λύση, $\mu$οντέλο δεδο$\mu$ένων και σχέδιο ανάπτυξης. Στην επιλογή της τεχνικής λύσης, λάβε υπόψη την αρχή «Simplicity First» — frontend $\mu$ε React αλλά χωρίς σύνθετη διαχείριση state όπως Redux· backend $\mu$ε Node.js αλλά χωρίς microservices· βάση δεδο$\mu$ένων SQLite αντί για PostgreSQL (στο στάδιο ανάπτυξης).

Το σχέδιο ανάπτυξης σπάει σε 5 βή$\mu$ατα: Step 1 ο σκελετός του backend API (server Express + SQLite + CRUD interfaces)· Step 2 οι βασικές σελίδες frontend (λίστα εγγραφών + φόρ$\mu$α προσθήκης)· Step 3 σύνδεση frontend-backend (κλήσεις API + ε$\mu$φάνιση δεδο$\mu$ένων)· Step 4 λειτουργίες κατηγοριών και φίλτρων· Step 5 $\mu$ηνιαία σύνοψη και απλά γραφή$\mu$ατα. Κάθε βή$\mu$α, όταν ολοκληρώνεται, αφήνει κάτι «που δουλεύει».

\booksubsection{Ρυθ$\mu$ίζεις το αρχείο $\mu$νή$\mu$ης του έργου}

Πριν ξεκινήσεις τον κώδικα, φτιάξε πρώτα το αρχείο $\mu$νή$\mu$ης του έργου. Είναι η βασική πρακτική της $\mu$ηχανικής συ$\mu$φραζο$\mu$ένων — να γνωρίζει η AI το υπόβαθρο του έργου στην αρχή κάθε συνο$\mu$ιλίας.

\begin{lstlisting}
# Έργο: full-stack εφαρ$\mu$ογή λογιστηρίου
## Επισκόπηση
Εφαρ$\mu$ογή προσωπικού λογιστηρίου, $\mu$ε καταγραφή, κατηγορίες και $\mu$ηνιαία σύνοψη.
## Τεχνολογίες
- Frontend: React + TypeScript + Vite
- Backend: Node.js + Express + TypeScript
- Βάση δεδο$\mu$ένων: SQLite ($\mu$ε better-sqlite3)
- Κοινοί ορισ$\mu$οί τύπων frontend-backend
## Συ$\mu$φωνίες API
- Βασική διαδρο$\mu$ή: /api
- Μορφή απάντησης: { success: boolean, data: any, error: string }
- Στυλ RESTful
## Μοντέλο δεδο$\mu$ένων
Record { id, amount, type, category, note, date }
## Κανόνες κώδικα
- TypeScript σε αυστηρή λειτουργία
- Εσοχή 2 κενών
- Συναρτήσεις $\mu$ε σχόλιο JSDoc
- Σχόλια και $\mu$ηνύ$\mu$ατα commit στα κινεζικά
\end{lstlisting}

\booksubsection{Υλοποιείς βή$\mu$α-βή$\mu$α}

Σύ$\mu$φωνα $\mu$ε το σχέδιο ανάπτυξης, κάθε βή$\mu$α ολοκληρώνεται σε $\mu$ια ξεχωριστή συνο$\mu$ιλία. Παρακάτω ένα παράδειγ$\mu$α εντολής για το πρώτο βή$\mu$α:

\begin{lstlisting}
Χρειάζο$\mu$αι τον σκελετό του backend API.
Απαιτήσεις:
1. Server Express, port 3001
2. Βάση SQLite, πίνακας records
3. Πεδία πίνακα records:
   id (TEXT PRIMARY KEY)
   amount (REAL)
   type (TEXT: 'income'|'expense')
   category (TEXT)
   note (TEXT)
   date (TEXT: YYYY-MM-DD)
4. API CRUD:
   GET /api/records - ανάκτηση όλων
   POST /api/records - προσθήκη
   PUT /api/records/:id - ενη$\mu$έρωση
   DELETE /api/records/:id - διαγραφή
5. Middleware διαχείρισης σφαλ$\mu$άτων
Δώσε πλήρη δο$\mu$ή κώδικα backend.
\end{lstlisting}

Πρόσεξε τα χαρακτηριστικά αυτής της εντολής: περιγράφει $\mu$όνο το τρέχον βή$\mu$α (τον σκελετό του backend), χωρίς ανά$\mu$ειξη απαιτήσεων frontend· απαριθ$\mu$εί ρητά τα πεδία και τα interfaces (δίνει στην AI επαρκείς περιορισ$\mu$ούς), ενώ η συγκεκρι$\mu$ένη υλοποίηση την αποφασίζει η AI· ακολουθεί την αρχή «Goal-Driven Execution» — λέει «τι θα γίνει» και όχι «πώς θα γίνει».

Μόλις η AI δώσει τον κώδικα, επαλήθευσε α$\mu$έσως: τρέξε τον server και δοκί$\mu$ασε $\mu$ε curl ή Postman αν κάθε interface δουλεύει σωστά. Όταν επιβεβαιώσεις ότι όλα είναι εντάξει, άνοιξε νέα συνο$\mu$ιλία για το Step 2 (βασικές σελίδες frontend). Κάθε βή$\mu$α ακολουθεί τον κύκλο «εντολή-επαλήθευση-νέα συνο$\mu$ιλία».

\booksection{Έργο 3: Εφαρ$\mu$ογή $\mu$ε ενίσχυση AI (προχωρη$\mu$ένο)}

\booksubsection{Σύλληψη του έργου}

Στόχος: $\mu$ια εφαρ$\mu$ογή που συνδέεται $\mu$ε το API $\mu$εγάλου $\mu$οντέλου — π.χ. ένας «βοηθός $\mu$ελέτης $\mu$ε AI»: ο χρήστης εισάγει το περιεχό$\mu$ενο που $\mu$ελετά και η AI φτιάχνει περίληψη, συνδέει σχετικές γνώσεις και προτείνει πηγές $\mu$άθησης. Το ιδιαίτερο του έργου: η ίδια η εφαρ$\mu$ογή σου καλεί AI, οπότε πρέπει να διαχειριστείς νέα ζητή$\mu$ατα όπως κλήσεις API, prompt engineering και διαχείριση σφαλ$\mu$άτων.

Το πιο ση$\mu$αντικό $\mu$άθη$\mu$α αυτού του έργου: όταν «χρησι$\mu$οποιείς AI για να γράψεις κώδικα που καλεί AI», πρέπει να προσέχεις ιδιαίτερα το prompt engineering. Τα prompt που χρησι$\mu$οποιεί η εφαρ$\mu$ογή σου όταν καλεί το API του $\mu$εγάλου $\mu$οντέλου πρέπει κι αυτά να ακολουθούν τις αρχές που συζητήσα$\mu$ε σε αυτόν τον οδηγό — σαφήνεια, συντο$\mu$ία, στόχευση.

\booksubsection{Κρίσι$\mu$α τεχνικά ση$\mu$εία}

Αυτό το έργο περιλα$\mu$βάνει $\mu$ερικά τεχνικά ση$\mu$εία που δεν συνάντησες στα δύο προηγού$\mu$ενα έργα:

\begin{itemize}
  \item Διαχείριση API κλειδιών: το κλειδί δεν $\mu$παίνει ποτέ απευθείας στον κώδικα — δούλεψέ το $\mu$ε $\mu$εταβλητές περιβάλλοντος.
  \item Διαχείριση ροής απαντήσεων: οι απαντήσεις των $\mu$εγάλων $\mu$οντέλων συνήθως έρχονται σε ροή και πρέπει να ε$\mu$φανίζονται σταδιακά.
  \item Έλεγχος κατανάλωσης Token: τα $\mu$εγάλα $\mu$οντέλα χρεώνονται ανά token, οπότε χρειάζεται έλεγχος της χρήσης σε επίπεδο εφαρ$\mu$ογής.
  \item Διαχείριση προτύπων Prompt: διαφορετικές λειτουργίες $\mu$ε διαφορετικά πρότυπα prompt, που χρειάζονται δο$\mu$η$\mu$ένη διαχείριση.
\end{itemize}

Για όλα αυτά τα τεχνικά ση$\mu$εία $\mu$πορείς να ζητήσεις τη βοήθεια της AI, αλλά πρέπει να καταλαβαίνεις το «γιατί» του καθενός — γιατί το κλειδί δεν $\mu$παίνει στον κώδικα; γιατί χρειάζεται ροή απαντήσεων; γιατί ελέγχεται η κατανάλωση token; Μόνο όταν καταλάβεις το «γιατί» $\mu$πορείς να κρίνεις αν η λύση της AI είναι λογική.

\booksection{Γενική λίστα πρακτικής διαδικασίας}

Ανεξαρτήτως $\mu$εγέθους έργου, ορίστε $\mu$ια γενική λίστα πρακτικής ροής για τον προγρα$\mu$$\mu$ατισ$\mu$ό $\mu$ε AI. Κάθε φορά που ξεκινάς νέο έργο, πέρνα τη λίστα $\mu$ία φορά για να βεβαιωθείς ότι δεν ξέχασες κάποιο κρίσι$\mu$ο βή$\mu$α.

\begin{center}
\begin{tabularx}{\linewidth}{|X|X|X|}
\hline
\textbf{Στάδιο} & \textbf{Τι κάνεις} & \textbf{Σχετική αρχή} \tabularnewline
\hline
1. Σκέψη απαιτήσεων & Πέρνα τα πέντε στοιχεία των απαιτήσεων και ξεκαθάρισε τι θα φτιάξεις & Think Before Coding \tabularnewline
2. Ανάλυση απαιτήσεων & Σπάς τη $\mu$εγάλη απαίτηση σε $\mu$ικρά βή$\mu$ατα που επαληθεύονται ανεξάρτητα & Simplicity First \tabularnewline
3. Γράφεις PRD & Το οργανώνεις σε δο$\mu$η$\mu$ένο έγγραφο για αναφορά & Goal-Driven \tabularnewline
4. Ρυθ$\mu$ίζεις αρχείο $\mu$νή$\mu$ης & Φτιάχνεις CLAUDE.md ή αντίστοιχο αρχείο & Harness Engineering \tabularnewline
5. Υλοποιείς βή$\mu$α-βή$\mu$α & Μία συνο$\mu$ιλία ανά βή$\mu$α: εντολή-επαλήθευση-επό$\mu$ενο & Surgical Changes \tabularnewline
6. Δοκι$\mu$ή και επαλήθευση & Α$\mu$έσως $\mu$ετά από κάθε βή$\mu$α δοκι$\mu$άζεις ότι η λειτουργία είναι σωστή & Iterate \& Verify \tabularnewline
7. Επιθεώρηση κώδικα & Έλεγχος λογικής της αρχιτεκτονικής, ασφάλειας, συντηρησι$\mu$ότητας & Iterate \& Verify \tabularnewline
8. Επανάληψη και βελτίωση & Βελτιώνεις σε κύκλους $\mu$ε βάση τα αποτελέσ$\mu$ατα δοκι$\mu$ών και επιθεώρησης & Iterate \& Verify \tabularnewline
\hline
\end{tabularx}
\end{center}

\begin{center}\small\color{gray}Γενική πρακτική λίστα ροής για τον προγρα$\mu$$\mu$ατισ$\mu$ό $\mu$ε AI\end{center}

\begin{tcbpractice}{Στην πράξη: ολοκλήρωσε το «Έργο 1» αυτή την εβδο$\mu$άδα}
Πάρε το «site προσωπικών εργαλείων» της ενότητας 1 του Κεφαλαίου 44 ως την πρώτη σου εργασία προγρα$\mu$$\mu$ατισ$\mu$ού $\mu$ε AI: ξόδεψε 5 λεπτά στα πέντε στοιχεία, στείλε την εντολή στην AI (οποιοδήποτε εργαλείο κάνει) και κάνε παραλαβή $\mu$ε την τριεπίπεδη επαλήθευση. Ολοκλήρωσέ το $\mu$έσα σε 30 λεπτά — και γεννήθηκε το πρώτο σου δη$\mu$ιούργη$\mu$α προγρα$\mu$$\mu$ατισ$\mu$ού $\mu$ε AI. Όταν τελειώσεις, φύλαξέ το καλά: είναι το ενθύ$\mu$ιο της $\mu$ετάβασής σου από «αναγνώστη» σε «πρακτικό».

\end{tcbpractice}

\begin{tcbsummary}{Σύνοψη κεφαλαίου}
\begin{itemize}
  \item Έργο 1 (για αρχή): καθαρά frontend site εργαλείων — ε$\mu$πειρία πλήρους ροής, ολοκληρώνεται σε 30 λεπτά
  \item Έργο 2 (ενδιά$\mu$εσο): full-stack σύστη$\mu$α λογιστικής — PRD + αρχείο $\mu$νή$\mu$ης + υλοποίηση σε 5 βή$\mu$ατα, κάθε βή$\mu$α «δουλεύει»
  \item Έργο 3 (προχωρη$\mu$ένο): βοηθός $\mu$ελέτης $\mu$ε AI — γράφεις $\mu$ε AI κώδικα που καλεί AI, $\mu$ε προσοχή στο prompt engineering
  \item Γενική ροή σε 8 βή$\mu$ατα: σκέψη απαιτήσεων ανάλυση PRD αρχείο $\mu$νή$\mu$ης σταδιακή υλοποίηση δοκι$\mu$ή επιθεώρηση επανάληψη
  \item Κάθε βή$\mu$α ακολουθεί τον κύκλο «εντολή-επαλήθευση-νέα συνο$\mu$ιλία»
\end{itemize}

\end{tcbsummary}
